\section{What to measure}
\label{sec:order-params}

A society's state is a set of weight vectors and an $N\times N$ matrix of
distrust. To compare states --- and, later, to compare a society of perceptrons
with a society of language models --- we need statistics that depend on neither
the dimension $K$ nor the parameterisation.

Two pairwise quantities do the work. The \emph{ideological overlap} of two agents
is the cosine of their weight vectors,
$\rho_{IJ}=\hat{w}_I\cdot\hat{w}_J/(\|\hat{w}_I\|\|\hat{w}_J\|)$, which is $+1$
when they agree on every issue and $-1$ when they disagree on every one. The
\emph{trust} a receiver $I$ places in an emitter $J$ is
$\eta_{J|I}=1-2\Phi(\hmu)$, which is $+1$ when $J$ is fully trusted and $-1$
when fully distrusted: the distrust field mapped onto a bounded, sign-carrying
scale. Writing $G_{IJ}=\kappa_I\kappa_J$ for the class indicator ($+1$ for
same-class pairs, $-1$ otherwise), the three pair correlations are
%
\begin{equation}
  \Rwmu = \bigl\langle (\eta_{I|J}+\eta_{J|I})\,\rho_{IJ} \bigr\rangle,
  \qquad
  \Rmuc = \bigl\langle G_{IJ}(\eta_{I|J}+\eta_{J|I}) \bigr\rangle,
  \qquad
  \Rcw  = \bigl\langle G_{IJ}\,\rho_{IJ} \bigr\rangle,
  \label{eq:correlations}
\end{equation}
%
averaged over unordered pairs and normalised to $[-1,1]$. $\Rwmu$ asks whether
agents trust those who agree with them; $\Rmuc$, the order parameter of
discrimination, asks whether trust follows class; $\Rcw$ asks whether opinion
follows class.

The pair correlations cannot distinguish a society split into two coherent camps
from one in which distrust is spread inconsistently, so we also measure social
balance in the sense of \citet{heider1946attitudes} and
\citet{cartwright1956structural}. A triple of agents is ideologically balanced
when $b^{I}_{IJK}=\rho_{IJ}\rho_{JK}\rho_{KI}>0$ and affectively balanced when
$b^{A}_{IJK}=\tfrac{1}{2}(\eta_{IJ}\eta_{JK}\eta_{KI}+\eta_{JI}\eta_{IK}\eta_{KJ})>0$;
the aggregates
%
\begin{equation}
  \BI = \bigl\langle b^{I}_{IJK}\bigr\rangle_{(IJK)},
  \qquad
  \BA = \bigl\langle b^{A}_{IJK}\bigr\rangle_{(IJK)}
  \label{eq:balance}
\end{equation}
%
run over all $\binom{N}{3}$ triples. A randomly initialised society has
$\BI\approx\BA\approx 0$: half its triples are frustrated. A society that has
settled into two internally coherent, mutually opposed blocs has
$\BI=\BA=1$, whatever the blocs are made of --- and a society that cannot settle
at all stays near zero or goes negative. Balance therefore separates
\emph{organised} disagreement from mere disorder, which is exactly the
distinction the phase diagram turns on. Both aggregates are computed in closed
form from matrix products rather than by enumerating triples
(Appendix~\ref{app:closed-form}).

Note that none of \eqref{eq:correlations} or \eqref{eq:balance} mentions weights,
gradients or $K$: they need only a signed opinion alignment and a signed trust per
pair, both of which are elicitable from a language agent. That is what makes the
phase diagram of \S\ref{sec:results} a prediction \S\ref{sec:llm} can test rather
than an analogy.
